\chapter{結論}

\section{本研究のまとめ}

本論文では、長期ロボットタスクにおける工程文脈依存的不確実性診断を対象とした。VLM を次行動生成器として用いるのではなく、現在画像、作業履歴、現在工程、次工程候補、工程条件リストを入力として、工程条件の不確実性を診断し、\texttt{act} / \texttt{reobserve} / \texttt{ask\_user} / \texttt{stop} / \texttt{defer} の処理分岐へ接続する方法を提案した。

本研究の課題は、工程文脈依存的不確実性の整理、VLM による構造化診断方法の設計、同一視覚状態・異なる工程文脈の対照評価、および実機を含む安全性検証である。本論文の貢献は、第一に、長期タスク中の不確実性を処理分岐の観点から整理した点、第二に、VLM 出力を証拠充足性、不確実性原因、工程関連性、処理分岐として構造化した点、第三に、文脈依存的な分岐切り替え能力を評価対象として明示した点にある。

\section{今後の課題}

今後は、より多様な家庭内タスクや異なるロボットプラットフォームへの適用、再観測戦略の自動化、ask\_user の対話設計、実行後の回復行動との統合が課題である。さらに、より大規模な計画器や world model と接続し、長期タスク全体の頑健性向上へ拡張する必要がある。
